- different kind of energy harvesters
    - piezo
    - solar
    - rotational
- State of knowledge
- variable reluctance energy harvester &&&&&&&&&
    - maximum theoretical power
- rectifying circuits proposed
    - simulations
    - real measurements
- choice of integrated circuits 
    - tests
- bluetooth module implementation





% Capitolo 2 — Background teorico
% Prima di parlare di circuiti, il lettore ha bisogno delle basi per capire le scelte progettuali. Copriresti:

% Faraday e il modello elettrico del VREH (EMF, resistenza e induttanza interna)
% Massimo trasferimento di potenza e impedance matching
% Principio di funzionamento dei rettificatori (FWR, VD, NVC)
% Principio dei convertitori DC-DC (Buck, Boost)

% Capitolo 3 — Topologie di Power Conditioning
% Qui presenti e confronti le topologie candidate per il tuo sistema a 6 fasi, con i loro vantaggi e svantaggi rispetto al tuo caso specifico.
% Capitolo 4 — Design e Simulazione
% Progettazione del circuito scelto, simulazioni SPICE o equivalente, analisi dei risultati.
% Capitolo 5 — Risultati e Validazione
% Confronto tra topologie, analisi di PCE e PEE ai diversi RPM.
% Capitolo 6 — Conclusioni